% Propietario conservado: /Users/ruben/Documents/New project/output/EXTREMA_MEDIA_RAZON_RECIPROCIDAD_ES_EN_20260918/ARTICULO/ES/source/libro/04_bilateral.tex
\chapter{Bilateral conservation and arithmetic-geometric transduction}
\label{x_reciprocidad:lib04:bilateral}

The holographic arithmetic-geometric coupling constant denotes
the generated register \(K\), with its ordered windows, phase origin,
orientation and incidences. Its production through the signed register and
Hadamard inversion is developed in \cref{x_reciprocidad:sec:k-registro}.
The rational reading \(\kappa\), the orbital frame \(O_K\) and the
operators transporting its realizations are different objects.
This distinction allows their functions to be composed: generation of the register,
arithmetic encoding, information transport and recovery of the
geometric readings.

The regional coordinates and \(\alpha\) are received from the
APP--TRIT--TPK genealogy, from the same enriched state and its joint discrete
structure of the continuum. The following constructions act on these
subsequent realizations. Their additional content is a dimension-independent
bilateral conservation, a reversible extension admitting
incoming memory and an archiving policy with uniformly
controlled recovery.

\section{Bilateral law for two isometric transports}
\label{x_reciprocidad:lib04:ley}

Let \(H,H'\) be Hilbert spaces and
\(\mathcal B,\mathcal C:H\to H'\) linear isometries, so that
\(\mathcal B^*\mathcal B=\mathcal C^*\mathcal C=I_H\).
The calligraphic typography distinguishes these transports from the arithmetic
base of the refinement. For \(a>0\), we define
\begin{equation}
 \begin{aligned}
 p&=a^2+a+1=a(a+1)+1,\\
 N&=(2a+1)p=(a+1)^3+a^3,\\
 Q_-&=a\mathcal B-\mathcal C,\qquad
 Q_+=(a+1)\mathcal B+\mathcal C.
 \end{aligned}
 \label{x_reciprocidad:lib04:def-q}
\end{equation}

\begin{theorem}[Balance of two isometric realizations]
\label{x_reciprocidad:lib04:teo-balance}
For the maps \eqref{x_reciprocidad:lib04:def-q},
\begin{equation}
 (a+1)Q_-^*Q_-+aQ_+^*Q_+=NI_H.
 \label{x_reciprocidad:lib04:balance}
\end{equation}
The identity holds in every dimension, including infinite Hilbert
spaces, and admits rectangular transports with common domain and
codomain.
\end{theorem}

\begin{proof}
Write \(X=\mathcal B^*\mathcal C+\mathcal C^*\mathcal B\).
The isometry of both transports gives
\begin{align*}
 Q_-^*Q_-&=(a^2+1)I-aX,\\
 Q_+^*Q_+&=((a+1)^2+1)I+(a+1)X.
\end{align*}
The weights \(a+1\) and \(a\) exactly cancel the terms
of \(X\). The remaining coefficient is
\[
 (a+1)(a^2+1)+a((a+1)^2+1)
 =2a^3+3a^2+3a+1=N.
\]
The compositions are of bounded operators and the calculation uses
only their products and adjoints. The proof therefore also applies
in infinite dimension.
\end{proof}

For \(a=8\), received from the preceding nonadic partition,
\begin{equation}
 9\|Q_-v\|^2+8\|Q_+v\|^2
 =(9^3+8^3)\|v\|^2=17\cdot73\,\|v\|^2.
 \label{x_reciprocidad:lib04:balance-canonico}
\end{equation}
The factor \(73=8\cdot9+1\) normalizes this general identity. In the
ternary module it is also realized as a determinant and similarity
factor; each of these functions retains its own hypotheses.
The balance \eqref{x_reciprocidad:lib04:balance} uses the isometries and admits both
transports of finite order and transports without that property.

\section{Factorization, recovery and compatibility of the channels}
\label{x_reciprocidad:lib04:factorizacion}

Define
\begin{equation}
 E_av=\frac1{\sqrt p}
 \begin{pmatrix}\sqrt{a(a+1)}\,\mathcal Bv\\\mathcal Cv\end{pmatrix},
 \qquad
 R_a=\frac1{\sqrt{2a+1}}
 \begin{pmatrix}\sqrt a&-\sqrt{a+1}\\
                 \sqrt{a+1}&\sqrt a\end{pmatrix}.
 \label{x_reciprocidad:lib04:factorizacion-er}
\end{equation}
The matrix \(R_a\) acts on the two channels and is understood to be tensored
with \(I_{H'}\). We check
\(E_a^*E_a=I_H\), \(R_a^*R_a=I_{H'\oplus H'}\), and
\begin{equation}
 W_av=R_aE_av
 =\frac1{\sqrt N}
 \begin{pmatrix}\sqrt{a+1}\,Q_-v\\\sqrt a\,Q_+v\end{pmatrix},
 \qquad W_a^*W_a=I_H.
 \label{x_reciprocidad:lib04:w}
\end{equation}
In the case \(a=8\), the first distribution has weights \(72/73\) and
\(1/73\); the second operation is a rotation between the channels.
This factorization expresses the role of \(72+1\) within the total
norm, distinguishing it from the linear restitution of
\eqref{x_reciprocidad:lib03:cadena-restauracion}.

\begin{proposition}[Bilateral recovery]
\label{x_reciprocidad:lib04:prop-recuperacion}
The unnormalized channels \(x=Q_-v\), \(y=Q_+v\) determine
\begin{equation}
 \mathcal Bv=\frac{x+y}{2a+1},\qquad
 \mathcal Cv=\frac{ay-(a+1)x}{2a+1},\qquad
 v=\mathcal B^*\frac{x+y}{2a+1}.
 \label{x_reciprocidad:lib04:inversa-canales}
\end{equation}
From the normalized channels we recover
\(v=W_a^*W_av\), with \(\|W_a^*\|=1\) when \(H\neq\{0\}\).
\end{proposition}

\begin{proof}
Adding the two definitions of \(Q_\pm\) eliminates \(\mathcal Cv\).
The combination \(aQ_+-(a+1)Q_-\) eliminates \(\mathcal Bv\).
The first isometry has left inverse \(\mathcal B^*\), which
proves the formulas. The last assertion follows from
\eqref{x_reciprocidad:lib04:w}: the adjoint of an isometry has norm one.
\end{proof}

The image \(\operatorname{im}W_a\) is closed and its orthogonal
projector is \(W_aW_a^*\). A normalized pair is compatible with a
unique input exactly when it belongs to this image; its orthogonal
component quantifies the compatibility defect.

If \(\mathcal B=I\) and \(\mathcal C=U\) is unitary, the two
\(Q_\pm\) are polynomials in the same operator. For the unnormalized
channels, the criterion simplifies to
\begin{equation}
 Q_+x=Q_-y.
 \label{x_reciprocidad:lib04:compatibilidad-unitaria}
\end{equation}
Indeed, the equality is equivalent to
\(U(x+y)=ay-(a+1)x\). Defining
\(v=(x+y)/(2a+1)\) yields by substitution
\(Q_-v=x\), \(Q_+v=y\). Necessity follows from the commutation
of the two polynomials. This simplified form belongs to the indicated unitary
domain; the image criterion covers general transports
between different spaces.

\section{Transformation of readers and relative symmetry}
\label{x_reciprocidad:lib04:lectores}

\begin{theorem}[Common reader in both channels]
\label{x_reciprocidad:lib04:teo-lector}
For every bounded self-adjoint operator \(F\) on \(H'\),
\begin{equation}
 W_a^*\operatorname{diag}(F,F)W_a
 =\frac{a(a+1)\mathcal B^*F\mathcal B+
                  \mathcal C^*F\mathcal C}{a(a+1)+1}.
 \label{x_reciprocidad:lib04:lector-general}
\end{equation}
If \(\mathcal B\) is unitary and \(\mathcal C=\mathcal BR\), with
\(R\) unitary, the expression, for \(G=\mathcal B^*F\mathcal B\),
is
\begin{equation}
 \mathcal T_a(G)=\frac{a(a+1)G+R^*GR}{a(a+1)+1}.
 \label{x_reciprocidad:lib04:lector-relativo}
\end{equation}
We have \(\mathcal T_a(G)=G\) if and only if \(GR=RG\).
\end{theorem}

\begin{proof}
The expansion of \((a+1)Q_-^*FQ_-+aQ_+^*FQ_+\) cancels the
terms \(\mathcal B^*F\mathcal C+\mathcal C^*F\mathcal B\).
The coefficients of the remaining terms are
\(a(a+1)(2a+1)\) for \(\mathcal B^*F\mathcal B\) and \(2a+1\)
for \(\mathcal C^*F\mathcal C\). Dividing by \(N\) proves
\eqref{x_reciprocidad:lib04:lector-general}. Substituting \(\mathcal C=\mathcal BR\)
produces \eqref{x_reciprocidad:lib04:lector-relativo}. Equality with \(G\) is equivalent
to \(R^*GR=G\), and the unitarity of \(R\) makes it equivalent to
\(GR=RG\).
\end{proof}

The reader \(F=I\) recovers the norm. For a general reader, the
identity determines its transformation, and the coincidence
\(\mathcal B^*F\mathcal B=\mathcal C^*F\mathcal C=G\) guarantees
its conservation in all states. In the canonical case,
\begin{equation}
 \mathcal T_8(G)=\frac{72G+R^*GR}{73}.
 \label{x_reciprocidad:lib04:lector-72}
\end{equation}
This formula relates conservation and symmetry through a specific
commutator. The observable that exactly transports an operator
\(A\) from the domain to the complete image is \(W_aAW_a^*\); its
ambient identity is \(W_aW_a^*\). The diagonal observable of
\eqref{x_reciprocidad:lib04:lector-general} is another choice, whose difference is
calculated by the same theorem.

\section{Reversible update with incoming memory}
\label{x_reciprocidad:lib04:reentrada}

Now let \(\mathcal B=I\), \(\mathcal C=U\), with \(U\)
unitary on the same space, and write
\begin{equation}
 A=\frac{\sqrt{a+1}}{\sqrt N}Q_-,\qquad
 D=\frac{\sqrt a}{\sqrt N}Q_+.
 \label{x_reciprocidad:lib04:ad}
\end{equation}
The operators \(A,D\) are normal and commute with one another and with their
adjoints, being polynomials in \(U\) and \(U^*\).

\begin{theorem}[Unitary extension]
\label{x_reciprocidad:lib04:teo-unitaria}
The update
\begin{equation}
 \mathscr R_a=
 \begin{pmatrix}A&-D^*\\D&A^*\end{pmatrix}:H\oplus H\to H\oplus H
 \label{x_reciprocidad:lib04:unitaria}
\end{equation}
is unitary. Its first column is \(W_a\), and an input
\((v,m)\) preserves \(\|v\|^2+\|m\|^2\).
\end{theorem}

\begin{proof}
The balance gives \(A^*A+D^*D=I\). Normality provides
\(AA^*+DD^*=I\). The cross products of
\(\mathscr R_a^*\mathscr R_a\) and
\(\mathscr R_a\mathscr R_a^*\) vanish by the preceding
commutations; their two diagonals are the identity. The first column
reproduces \eqref{x_reciprocidad:lib04:w}, and conservation of the joint norm
is the unitary identity applied to \((v,m)\).
\end{proof}

For data \(z=(z_1,z_2)\), the explicit inversion is
\begin{equation}
 v_{\mathrm{rec}}=A^*z_1+D^*z_2,\qquad
 m_{\mathrm{rec}}=-Dz_1+Az_2.
 \label{x_reciprocidad:lib04:inversa-unitaria}
\end{equation}
The hypothesis of zero incoming memory is checked through
\(m_{\mathrm{rec}}=0\). If \(e\) is a perturbation of the data,
unitarity gives
\begin{equation}
 \|e\|^2=\|e_{\mathrm{rec}}\|^2+\|e_{\mathrm{mem}}\|^2.
 \label{x_reciprocidad:lib04:descomposicion-ruido}
\end{equation}
The second component detects the part incompatible with that input.
An error belonging to \(\operatorname{im}W_a\) has zero memory
residual and alters the recovered input. Correction of the integer
register also uses its discrete separation distances.


% BEGIN OWNER /Users/ruben/Documents/New project/output/EXTREMA_MEDIA_RAZON_RECIPROCIDAD_ES_EN_20260918/ARTICULO/ES/source/libro/memoria_reducida_20260916.tex
\section{Causal reduction of the bilateral update}
\label{x_reciprocidad:lib04:reduccion-causal}

The preceding unitary update realizes the transport of the state and its
memory in a common space. Observation of the first component admits an
equivalent description that explicitly retains the contribution of the
record. This description follows from the transport already constructed,
with the same coefficients and orientation: memory removed from the list
of variables reappears as dependence on earlier states.

Fix the operators \(A,D\) of \eqref{x_reciprocidad:lib04:ad} and a realization
\(H\oplus H\). For \(n\geq0\), let
\begin{equation}
 v_{n+1}=Av_n-D^*m_n,\qquad
 m_{n+1}=Dv_n+A^*m_n.
 \label{x_reciprocidad:lib04:recurrencia-memoria}
\end{equation}
The joint update is \(\mathscr R_a\) and preserves
\(\|v_n\|^2+\|m_n\|^2\).

\begin{proposition}[Exact elimination of the auxiliary record]
\label{x_reciprocidad:lib04:prop-memoria-reducida}
For every \(n\geq0\), interpreting the empty sum as zero,
\begin{align}
 m_n&=(A^*)^n m_0+
       \sum_{k=0}^{n-1}(A^*)^{n-1-k}Dv_k,
       \label{x_reciprocidad:lib04:memoria-explicita}\\
 v_{n+1}&=Av_n-D^*(A^*)^n m_0
       +\sum_{k=0}^{n-1}\mathscr K_{n-1-k}v_k,
       \label{x_reciprocidad:lib04:evolucion-reducida}\\
 \mathscr K_j&=-D^*(A^*)^jD\quad(j\geq0).
       \label{x_reciprocidad:lib04:nucleo-retardo}
\end{align}
The reduced equation, together with \eqref{x_reciprocidad:lib04:memoria-explicita}, is
equivalent to the joint recurrence for the same data \((v_0,m_0)\).
\end{proposition}
\begin{proof}
The first identity holds for \(n=0\). If it holds for \(n\), the second
equation in \eqref{x_reciprocidad:lib04:recurrencia-memoria} gives
\[
 m_{n+1}=(A^*)^{n+1}m_0+
 \sum_{k=0}^{n-1}(A^*)^{n-k}Dv_k+Dv_n,
\]
which is the identity at \(n+1\). Substitution at depth \(n\) into the
first equation proves \eqref{x_reciprocidad:lib04:evolucion-reducida}. Conversely, if
\(v_n\) satisfies the reduced equation and \(m_n\) is defined by
\eqref{x_reciprocidad:lib04:memoria-explicita}, the same difference of sums proves
\(m_{n+1}=Dv_n+A^*m_n\); substitution recovers the other equation.
\end{proof}

The kernel \(\mathscr K_j\) describes a transfer into the record,
\(j\) updates within it, and re-entry into the observed component.
The term \(-D^*(A^*)^nm_0\) preserves the initial memory. Omitting it
corresponds to the particular case \(m_0=0\), rather than an identity of
the general transport. Equation \eqref{x_reciprocidad:lib04:evolucion-reducida} involves
only the current and earlier states. The causal character of this reduction
follows from the order of the joint recurrence.

The normalization of the balance also provides an explicit estimate.
Since \(Q_-=aI-U\), we have
\(\|A\|^2\leq r_a=(a+1)^3/((a+1)^3+a^3)<1\), whereas
\(D^*D\leq I\). Consequently,
\begin{equation}
 \|\mathscr K_j\|\leq r_a^{j/2},\qquad
 \sum_{j\geq0}\|\mathscr K_j\|
 \leq\frac{1}{1-\sqrt{r_a}}.
 \label{x_reciprocidad:lib04:cota-nucleo-retardo}
\end{equation}
This bound controls the weight of memory in the reduced equation; the norm
of the joint state remains conserved by the unitary update.

The proposition is the temporal realization of the Schur reduction in
\cref{x_reciprocidad:mem:prop-schur}. For the formal series
\(V(z)=\sum_{n\geq0}v_nz^n\) and
\(M(z)=\sum_{n\geq0}m_nz^n\), the recurrences yield
\(V=v_0+zAV-zD^*M\) and \(M=m_0+zDV+zA^*M\). Eliminating \(M\),
\begin{equation}
 V(z)=\bigl[I-zA+z^2D^*(I-zA^*)^{-1}D\bigr]^{-1}
       \bigl[v_0-zD^*(I-zA^*)^{-1}m_0\bigr].
 \label{x_reciprocidad:lib04:schur-bilateral}
\end{equation}
The formal inverses exist because their constant terms are identities;
the state series and the joint resolvent also converge for \(|z|<1\).
The initial source and the effective operator transform together. A readout
that also acts on \(m_n\) retains its contribution through
\eqref{x_reciprocidad:lib04:memoria-explicita}, in accordance with \eqref{x_reciprocidad:mem:schur-lector}.

Re-entry here updates an incoming record. By contrast, the archiving policy
below retains each complement in a distinct coordinate. Each operation has
its own conservation law. The kernel \(\mathscr K_j\), Green propagators,
and a dimensional force have different types: the first eliminates a state
component, Green propagators solve source equations with support conditions, and
the force adds a dynamical realization and its units.

% END OWNER


\section{Nonstationary register and uniform convergence of the terminal component to zero}
\label{x_reciprocidad:lib04:archivo}

At each stage let \(H_n,H_{n+1}\) be Hilbert spaces and
\(\mathcal B_n,\mathcal C_n:H_n\to H_{n+1}\) isometries. The
parameter \(a>0\) remains fixed. Let \(A_n,D_n\) be the normalized
blocks of \eqref{x_reciprocidad:lib04:w}; we continue through the first and
record the second:
\begin{equation}
 R_0=I,\qquad R_{n+1}=A_nR_n,\qquad m_n=D_nR_nv.
 \label{x_reciprocidad:lib04:politica-archivo}
\end{equation}
This policy has a domain distinct from the reintroduction of
memory in \eqref{x_reciprocidad:lib04:unitaria}; each preserves its update.

\begin{lemma}[Uniform transfer]
\label{x_reciprocidad:lib04:lem-tasa}
At every stage,
\begin{equation}
 \|A_n\|^2\leq r_a:={\frac{(a+1)^3}{(a+1)^3+a^3}}<1,
 \quad D_n^*D_n\geq(1-r_a)I,
 \quad\|R_n\|^2\leq r_a^n.
 \label{x_reciprocidad:lib04:tasa}
\end{equation}
\end{lemma}

\begin{proof}
The triangle inequality gives
\(\|a\mathcal B_n-\mathcal C_n\|\leq a+1\). Its normalization
provides the first bound. The difference
\(N-(a+1)^3=a^3>0\) proves \(r_a<1\). The local balance implies
\(D_n^*D_n=I-A_n^*A_n\), and multiplying the bounds of the continued
channels produces the estimate of \(R_n\).
\end{proof}

\begin{theorem}[Complete memory at arbitrary depth]
\label{x_reciprocidad:lib04:teo-archivo}
The finite map
\begin{equation}
 Z_Nv=(R_Nv,D_0v,D_1R_1v,\ldots,D_{N-1}R_{N-1}v)
 \label{x_reciprocidad:lib04:zn}
\end{equation}
is isometric. The infinite archive
\begin{equation}
 Z_\infty v=(D_nR_nv)_{n\geq0}
 \label{x_reciprocidad:lib04:zinfty}
\end{equation}
belongs to the Hilbert direct sum of the output spaces and satisfies
\begin{equation}
 \|v\|^2=\sum_{n\geq0}\|m_n\|^2,\qquad
 v=\sum_{n\geq0}R_n^*D_n^*m_n,
 \qquad Z_\infty^*Z_\infty=I.
 \label{x_reciprocidad:lib04:recuperacion-infinita}
\end{equation}
The inversion series converges in norm and the adjoint has norm one.
\end{theorem}

\begin{proof}
The identity \(A_n^*A_n+D_n^*D_n=I\), conjugated by \(R_n\),
produces
\[
 R_n^*R_n-R_{n+1}^*R_{n+1}=R_n^*D_n^*D_nR_n.
\]
Summing from zero to \(N-1\) gives
\begin{equation}
 I=R_N^*R_N+\sum_{n<N}R_n^*D_n^*D_nR_n.
 \label{x_reciprocidad:lib04:telescopia}
\end{equation}
This is the Gram identity of \(Z_N\). By
\eqref{x_reciprocidad:lib04:tasa}, the first term tends to zero in operator
norm. The limit identity proves the quadratic convergence of the
archive and its isometry. The adjoint of that isometry is bounded; the
truncations of every square-summable archive converge in the Hilbert
direct sum, and their image under the adjoint proves convergence of the
series in \eqref{x_reciprocidad:lib04:recuperacion-infinita}. Applying it to the exact
data recovers \(v\).
\end{proof}

With only the first \(N\) complements, the truncated adjoint has
exact residual
\begin{equation}
 v-\sum_{n<N}R_n^*D_n^*m_n=R_N^*R_Nv,
 \qquad
 \left\|v-\sum_{n<N}R_n^*D_n^*m_n\right\|
 \leq r_a^N\|v\|.
 \label{x_reciprocidad:lib04:sesgo-truncado}
\end{equation}
If the joint data error is at most \(\varepsilon\), the
total bound is \(\varepsilon+r_a^N\|v\|\). With \(Z_N\), including
its terminal, or with \(Z_\infty\), the adjoint recovers with error at
most \(\varepsilon\).

The preceding policy exhausts the terminal through the proved uniform bound.
The preceding evolutions that have a positive terminal
\(A_\infty\neq0\) preserve in their archive the component
\(A_\infty^{1/2}v\). Arbitrary depth and removal of the
terminal are therefore different conclusions and retain their
respective hypotheses.

\begin{proposition}[Inverse of finite memories without a terminal]
\label{x_reciprocidad:lib04:prop-inversa-finita}
For \(N\geq1\), let \(M_Nv=(D_nR_nv)_{n<N}\). Then
\begin{equation}
 M_N^*M_N=I-R_N^*R_N\geq(1-r_a^N)I,
 \label{x_reciprocidad:lib04:gram-finito}
\end{equation}
and the least-squares inverse
\begin{equation}
 L_N=(I-R_N^*R_N)^{-1}M_N^*,\qquad
 L_NM_N=I,\qquad
 \|L_N\|\leq(1-r_a^N)^{-1/2}
 \label{x_reciprocidad:lib04:inversa-finita}
\end{equation}
recovers the input exactly with exact data.
\end{proposition}

\begin{proof}
The Gram identity is \eqref{x_reciprocidad:lib04:telescopia}; the bound follows from
\eqref{x_reciprocidad:lib04:tasa}. The positive operator is invertible by its strict
lower bound. Composition proves \(L_NM_N=I\), and
\(L_NL_N^*=(M_N^*M_N)^{-1}\) proves the stated norm.
\end{proof}

This finite inverse removes the bias of
\eqref{x_reciprocidad:lib04:sesgo-truncado}; its noise amplification is
specified by \eqref{x_reciprocidad:lib04:inversa-finita}. The truncated adjoint,
the finite inverse and the adjoint of the complete archive are thus three
operators with distinct behaviors.

For \(a=8\), the general bound is \(r_a=729/1241\). In the case
\(\mathcal B=I\), \(\mathcal C=S^4\), with \(S\) the cycle of
twelve positions, the fixed projector
\(P_0=(I+S^4+S^8)/3\) gives
\begin{equation}
 A^*A=\frac{a+1}{(2a+1)p}
       \bigl((a-1)^2P_0+p(I-P_0)\bigr).
 \label{x_reciprocidad:lib04:gram-tasa-ternaria}
\end{equation}
Since \(p-(a-1)^2=3a>0\), its norm is
\((a+1)/(2a+1)\). At \(a=8\) we obtain \(9/17\); the eigenvalue
of the fixed sector is \(441/1241\). These rates belong to the
declared transports and policy.

If both channels are refined in a tree, each complete level is also
isometric by composition. Summing the norms of all complete layers
would repeatedly count the same initial norm. Conservation
is formulated through the refinement maps between levels or
through the archive of complements in \eqref{x_reciprocidad:lib04:zinfty}.

\section{Geometric forms and conservation of their transport}
\label{x_reciprocidad:lib04:formas}

Let \(b,b'\) be nondegenerate forms, symmetric or alternating, on
the corresponding spaces, and suppose
\(\mathcal B^*b'\mathcal B=\mathcal C^*b'\mathcal C=b\).
The same bilinear expansion proves
\begin{equation}
 (a+1)Q_-^*b'Q_-+aQ_+^*b'Q_+=Nb.
 \label{x_reciprocidad:lib04:balance-formas}
\end{equation}
Indeed, the mixed terms
\(\mathcal B^*b'\mathcal C+\mathcal C^*b'\mathcal B\) cancel
with the same weights; the pure terms sum to \(Nb\).
The transported form is thus preserved without requiring positivity.
The contraction and noise bounds of the preceding subsection use,
by contrast, positive Hilbert norms. A Lorentzian or
alternating form preserves \eqref{x_reciprocidad:lib04:balance-formas}, with its own
geometric interpretations, and does not replace positivity in those
inequalities.

\section{Channel correlation and inverse reading problem}
\label{x_reciprocidad:lib04:correlacion}

For \(v\neq0\), define
\begin{equation}
 \chi(v)=\frac{\operatorname{Re}\langle\mathcal Bv,\mathcal Cv\rangle}
                      {\|v\|^2},\qquad -1\leq\chi(v)\leq1.
 \label{x_reciprocidad:lib04:chi}
\end{equation}
The bound is Cauchy--Schwarz applied to the two isometric images.
The normalized readings are
\begin{equation}
 s_-=a^2+1-2a\chi,\qquad
 s_+=(a+1)^2+1+2(a+1)\chi.
 \label{x_reciprocidad:lib04:s-pm}
\end{equation}
For \(a=8\),
\begin{equation}
 s_-=65-16\chi,\qquad s_+=82+18\chi,\qquad
 9s_-+8s_+=1241.
 \label{x_reciprocidad:lib04:s-canonico}
\end{equation}
The two individual norms contain a common real correlation. The
balance is an exact relation between them, not two independent arithmetic
determinations of a constant.

Using the prior output \(\alpha_A\) in the inverse test
\(s_-=10^4\alpha_A\) determines the complementary reading
\begin{equation}
 s_+=73+\frac98(73-10^4\alpha_A)
       \simeq73{,}029783595557.
 \label{x_reciprocidad:lib04:ensayo-alpha}
\end{equation}
The equality is obtained by substitution into
\eqref{x_reciprocidad:lib04:s-canonico}. The weighted average
\((9s_-+8s_+)/17\) remains exactly equal to \(73\).
The comparison preserves the prior genealogy of \(\alpha_A\): it is
a condition of the inverse test, not another independent producer of
that output. A unitary change can alter \(\chi\) and the two
readings while maintaining the balance, even if a particular relation
of ternary order changes.

\section{Cyclic frame and transport of responses}
\label{x_reciprocidad:lib04:marco}

On the carrier of the generated register, let
\(H=\mathbb R^{12}\), \((Sx)_i=x_{i+1}\), with cyclic indices,
and
\begin{equation}
 O_K=(K,SK,\ldots,S^{11}K),\qquad
 589609I\leq O_K^*O_K\leq39225169I.
 \label{x_reciprocidad:lib04:gram-k}
\end{equation}
The bounds are those of the cyclic frame of the canonical register,
proved in \cref{x_reciprocidad:lib01:marco-ciclico}.
Let \(Z\) be one of the complete isometries
\(Z_N,Z_\infty\), or a composition with the preceding archives
that preserves its terminals and complements.

\begin{theorem}[Preservation of frame conditioning]
\label{x_reciprocidad:lib04:teo-marco}
If \(x=O_K\beta\), the datum \(z=Zx\) determines
\begin{equation}
 \beta=O_K^{-1}Z^*z,\qquad
 \|\delta\beta\|\leq\frac{\|\delta z\|}{\sqrt{589609}}.
 \label{x_reciprocidad:lib04:inversa-marco}
\end{equation}
The twelve columns of \(ZO_K\) form a basis of its image, even
when the archive has infinitely many coordinates.
\end{theorem}

\begin{proof}
The equality \(Z^*Z=I\) gives
\((ZO_K)^*(ZO_K)=O_K^*O_K\). The Gram bounds are preserved and
\(O_K\) is invertible by the strict lower bound. Composing their
inverses proves recovery, and the norm
\(\|O_K^{-1}\|\leq589609^{-1/2}\), together with
\(\|Z^*\|=1\), gives stability. The number of columns and their
independence determine the dimension of the image, without identifying it
with the entire storage space.
\end{proof}

For an operator \(H_0\) commuting with \(S\), the complete
vector response \(y=H_0K\) determines
\begin{equation}
 H_0=O_yO_K^{-1}.
 \label{x_reciprocidad:lib04:identificacion-filtro}
\end{equation}
Indeed, commutation implies
\(H_0O_K=(y,Sy,\ldots,S^{11}y)=O_y\). If only
\(z=Zy+e\) is received, define \(\widehat y=Z^*z\) and
\(\widehat H_0=O_{\widehat y}O_K^{-1}\). Then
\begin{equation}
 \|\widehat H_0-H_0\|
 \leq\sqrt{\frac{12}{589609}}\,\|e\|.
 \label{x_reciprocidad:lib04:error-filtro}
\end{equation}
The orbital matrix of \(\delta y\) has twelve columns of norm
\(\|\delta y\|\), so its operator norm does not exceed
\(\sqrt{12}\|\delta y\|\). This and
\eqref{x_reciprocidad:lib04:gram-k} prove \eqref{x_reciprocidad:lib04:error-filtro}.
For a general operator the protocol requires the twelve responses
\(H_0S^jK\); a response in the cyclic case means its twelve
components, not a scalar measurement.

In the transported image, the shift is
\(S_Z=ZSZ^*\), and satisfies \(S_Z^jZ=ZS^j\). This equality
determines the phase advance in the archive; a translation of cells
in the ambient space represents that advance only when it coincides with
the transported operator.

\section{Positional reading and scalar precision}
\label{x_reciprocidad:lib04:kappa}

Fix \(B_0=1000\), \(d=B_0^{12}-1\), and the row
\begin{equation}
 \ell=\frac{(B_0^{11},B_0^{10},\ldots,1)}{d}.
 \label{x_reciprocidad:lib04:fila-posicional}
\end{equation}
Its linear extension acts on \(H\), and in the word domain
\(\kappa(K)=\ell K\). The representative of the transported reader is
\(Z\ell^*\), so that
\begin{equation}
 \kappa(K)=\langle Z\ell^*,ZK\rangle,\qquad
 |\delta\kappa|\leq\|\ell\|\,\|e\|.
 \label{x_reciprocidad:lib04:lectura-kappa}
\end{equation}
The geometric sum provides
\begin{equation}
 \|\ell\|^2
 =\frac{B_0^{12}+1}{(B_0^2-1)(B_0^{12}-1)}.
 \label{x_reciprocidad:lib04:norma-fila}
\end{equation}
Indeed, the numerator of the sum of squares is
\((B_0^{24}-1)/(B_0^2-1)\), which divided by
\((B_0^{12}-1)^2\) gives \eqref{x_reciprocidad:lib04:norma-fila}.

With only \(N\) complements and the truncated adjoint, an additional
error of at most \(\|\ell\|r_a^N\|K\|\) appears. The finite inverse
\eqref{x_reciprocidad:lib04:inversa-finita} removes that bias and retains its explicit
noise amplification bound.

In the domain \(\{0,\ldots,999\}^{12}\), the exact reading
recovers \(N_K=d\kappa\) and its twelve blocks through successive Euclidean
divisions. If
\begin{equation}
 |\widetilde\kappa-\kappa|<\frac1{2d},
 \label{x_reciprocidad:lib04:radio-escalar}
\end{equation}
rounding \(d\widetilde\kappa\) to the nearest integer
recovers \(N_K\): its distance to the true integer is less than
half a unit. The base, length, origin and orientation accompany
that representation. The scalar threshold of
\eqref{x_reciprocidad:lib04:radio-escalar} is distinct from the vector bounds
of the memory registers.

\section{Incidence, charge and exact restitution of the register}
\label{x_reciprocidad:lib04:incidencia}

Let \(\mathbf1=(1,\ldots,1)^T\) and
\(\Pi=I-\mathbf1\mathbf1^*/12\). The incidence isometry of
\cref{x_reciprocidad:nuc05:isometria} is
\begin{equation}
 Fv=(\mu(v),\mathcal I\Pi v/6),\qquad
 \mu(v)=\frac{\mathbf1^*v}{12},\qquad
 \|(\mu,y)\|_Y^2=12\mu^2+\|y\|^2.
 \label{x_reciprocidad:lib04:f-incidencia}
\end{equation}
Its codomain is its declared image. The Gram matrix
\(\mathcal I^*\mathcal I=36I+30\mathbf1\mathbf1^*\) gives
\(\|\mathcal I\Pi v\|^2=36\|\Pi v\|^2\), and
\(12\mu(v)^2+\|\Pi v\|^2=\|v\|^2\), whence \(F^*F=I\).
Thus, from \(z=Zv+e\),
\begin{equation}
 \widehat{Fv}=FZ^*z,\qquad \|FZ^*e\|_Y\leq\|e\|.
 \label{x_reciprocidad:lib04:error-incidencia}
\end{equation}
The unnormalized incidence \(\mathcal I\Pi v\) has error at
most \(6\|e\|\); the signed register \(H_{12}v\), whose
normalization satisfies \(H_{12}^*H_{12}=4I\), has error at most
\(2\|e\|\).

The direction \(u=P_3\Pi K\) and the projector
\(P_{10}=\Pi-uu^*/\|u\|^2\), with \(u\neq0\), are recovered
with their prior readers and chart. The absolute selection
\(\{j:K_j\geq729\}\) has a discontinuity at \(K_4=729\).
Integer restitution is carried out before applying that threshold.

\begin{proposition}[Decoding after the complete archive]
\label{x_reciprocidad:lib04:prop-redondeo}
Let \(K\) be integer and \(z=ZK+e\), with \(Z^*Z=I\). If
\(\|e\|<1/2\), rounding each coordinate of \(Z^*z\) recovers
\(K\) exactly. Composition restores \(\kappa\), the
signatures and the geometric readers defined on that register.
\end{proposition}

\begin{proof}
The contractivity of \(Z^*\) gives
\(\|Z^*z-K\|\leq\|e\|<1/2\). Each coordinate error is
bounded by the joint norm, so the unique integer coordinate
at a distance less than half a unit is the original one. Once
\(K\) is recovered, evaluating any of its exact readers
restores its value.
\end{proof}

In particular, the support \(\{4,6,9,10\}\) is recovered, including
its boundary component. The other marks involved in an
exceptional flag are preserved through their own channels:
restitution of \(K\) exactly corrects the information whose
sufficient arguments lie in it.

\section{Preservation of the conditioning of two memories}
\label{x_reciprocidad:lib04:dos-memorias}

Retain the preceding readers
\begin{equation}
 T_j=\frac{8I+S^j}{9},\qquad
 D_j=\frac{\sqrt8}{9}(S^j-I),\qquad
 Mv=(D_3v,D_4T_3v).
 \label{x_reciprocidad:lib04:m-antecedente}
\end{equation}
The \(D_j\) in this equation are the preceding complements;
the second bilateral block \(D\) of \eqref{x_reciprocidad:lib04:ad} retains
its own definition. We have \(\ker M=\mathbb R\mathbf1\)
and an inverse \(L_M\) of norm \(9/4\) on the centered sector,
as proved in \cref{x_reciprocidad:nuclear:11}.

Let \(E=W_a\oplus W_a\) be in the space of the two memories.
Its isometry implies
\begin{equation}
 (EM)^*(EM)=M^*M,\qquad
 \ker(EM)=\mathbb R\mathbf1.
 \label{x_reciprocidad:lib04:gram-m-preservado}
\end{equation}
For \(z=EMv+e\) and exact charge
\(q=\mathbf1^*v\), the reconstruction is
\begin{equation}
 \widehat v=\frac q{12}\mathbf1+L_ME^*z,\qquad
 \|\widehat v-v\|\leq\frac94\|e\|.
 \label{x_reciprocidad:lib04:inversa-m}
\end{equation}
The proof consists in applying \(E^*E=I\), using
\(L_MM=\Pi\) and \(\|E^*\|=1\). If the charge has error
\(\delta q\), the uniform and centered sectors are orthogonal and
\begin{equation}
 \|\widehat v-v\|^2
 \leq\frac{|\delta q|^2}{12}+\frac{81}{16}\|e\|^2.
 \label{x_reciprocidad:lib04:error-carga}
\end{equation}

The distances of all words in the integer image are preserved
exactly by \(E\). Therefore, the restitution radii of
\cref{x_reciprocidad:lib02:radios-memoria} remain
\begin{equation}
 r_0=\frac{2\sqrt{146}}{81}\quad\hbox{uncharged},\qquad
 r_q=\frac{4\sqrt{73}}{81}\quad\hbox{with exact charge}.
 \label{x_reciprocidad:lib04:radios}
\end{equation}
Within the first radius, \(\Pi K\) is recovered, and with it the
direction and its projector when they are defined. Within the second,
preserving \(q\), we recover \(K\) and the absolute threshold
selector. Subsequently composing another complete isometry, finite or infinite,
maintains the same guarantees: its adjoint is applied before the
decoder. The radius is always expressed in the joint norm
of the data with the fixed normalization.

This composition redistributes the information of \(M\) without degrading
its conditioning and without creating the uniform channel that \(M\) removes.
Gram conservation distinguishes this property from a numerical
improvement obtained by changing the reader or the noise scale.

\section{Orbital norm and polar normalization of the frame}
\label{x_reciprocidad:lib04:normalizacion}

The preceding inversion satisfies
\(K=H_{12}U_{\mathrm{sgn}}/4\),
\(H_{12}^*H_{12}=4I\). Therefore
\begin{equation}
 \|K\|^2=\frac14\|U_{\mathrm{sgn}}\|^2=4251257,
 \qquad \|U_{\mathrm{sgn}}\|^2=17005028.
 \label{x_reciprocidad:lib04:norma-k}
\end{equation}
Each translation \(S^j\) is orthogonal and preserves
\(\nu=\sqrt{4251257}\). Hence
\begin{equation}
 \operatorname{tr}(O_K^*O_K)=12\nu^2=51015084.
 \label{x_reciprocidad:lib04:traza-marco}
\end{equation}
Division by \(\nu\) normalizes each orbital column to norm
one and preserves its order. It is uniform over the orbit and under its
isometric transports. Different registers retain their own
norms, and the normalized columns continue to have a nonscalar
Gram matrix: its eigenvalues are the preceding ones divided by
\(4251257\).

Operator normalization uses the complete Gram matrix:
\begin{equation}
 G_K=O_K^*O_K,\qquad
 F_K=O_KG_K^{-1/2},\qquad
 F_K^*F_K=I.
 \label{x_reciprocidad:lib04:polar}
\end{equation}
The inverse root exists by \eqref{x_reciprocidad:lib04:gram-k}. The equality is
checked by multiplying the Gram matrix by \(G_K^{-1/2}\) on both sides.
Thus \(ZF_K\) is isometric at any depth.
The full recovery of the frame is
\begin{equation}
 O_K=F_KG_K^{1/2}.
 \label{x_reciprocidad:lib04:polar-inversa}
\end{equation}
The pair \((F_K,G_K)\) preserves the complete information.
The unitary factor preserves an oriented orthonormal frame in the
fixed charts, while the amplitudes also reside in
\(G_K\). Multiplying \(K\) by a positive scalar preserves
\(F_K\) and changes \(G_K\); \(K\) and \(-K\) share a
Gram matrix and have opposite polar factors, with different linear reading
\(\kappa\). Each example shows why an isolated factor does not
replace the pair.

Translation of the origin preserves \(\nu\), but in the fixed positional
chart
\begin{equation}
 \kappa(SK)=1000\kappa(K)-K_1.
 \label{x_reciprocidad:lib04:cambio-origen}
\end{equation}
This is obtained by shifting the positional numerator and using
\(d=1000^{12}-1\). To preserve the same reading under a chart
change, the row \(\ell\) is also transported by the inverse
of that change. A general unitary chart may leave the lattice of
integer blocks; decoding by Euclidean divisions is performed
in the recovered block chart.

\section{Exact example of transduction and error correction}
\label{x_reciprocidad:lib04:ejemplo}

Take \(a=8\), \(U=S^4\) and the canonical register
\begin{equation}
 \begin{split}
 K=(&234,543,140,729,659,824,\\
    &621,58,914,794,146,601).
 \end{split}
 \label{x_reciprocidad:lib04:k-ejemplo}
\end{equation}
Denoting by \((v)_{1:6}\) and \((v)_{7:12}\) its two blocks of
six coordinates, the unnormalized readings are
\begin{align}
 (Q_-K)_{1:6}&=(1213,3520,499,5774,4358,5798),\nonumber\\
 (Q_-K)_{7:12}&=(4822,-137,7078,5809,1028,4079),\nonumber\\
 (Q_+K)_{1:6}&=(2765,5711,1881,6619,6845,8210),\nonumber\\
 (Q_+K)_{7:12}&=(5735,1123,8460,7689,1454,6138).
 \label{x_reciprocidad:lib04:canales-ejemplo}
\end{align}
Their sum, divided by \(17\), recovers all blocks of \(K\).
The matrix identity
\(9Q_-^*Q_-+8Q_+^*Q_+=1241I\) precedes its evaluation on the
specific vector.

To introduce a controlled error, the two memories are rationalized
as \(z_0=MK/\sqrt8\), and \(1/10\) is added to the first
coordinate. Each block of twelve components is then transported
through the two bilateral channels. The inverse by summation recovers the
perturbed memory data, and the charge decoder receives
\(q=6263\) together with these data, without receiving the sought register.
The error norm in the original memories is
\(\sqrt8/10<r_q\). Isometric normalization preserves that norm;
the unnormalized channels merely facilitate exact rational calculation.

In this instance, the candidates of the floor-and-ceiling reconstruction
satisfying the charge are \(924\). The minimum rationalized
squared distance is \(1/100\), and the decoder recovers
\eqref{x_reciprocidad:lib04:k-ejemplo}. Uniqueness follows from the strict radius
\eqref{x_reciprocidad:lib04:radios}, not from the mere size of the enumeration.
The recovered arithmetic reading is
\begin{equation}
 \kappa=
 \frac{234543140729659824621058914794146601}
      {999999999999999999999999999999999999},
 \label{x_reciprocidad:lib04:kappa-ejemplo}
\end{equation}
and the incidence selection is again \(\{4,6,9,10\}\), with
the same threshold \(729\). The example introduces noise before
transport. The proofs of contractivity of the adjoint also cover
arbitrary subsequent errors whose joint norm lies within the
declared radius.

The combined operational capacity is to transport and recover a generated
register, maintain the conditioning of its responses and restore
its arithmetic and incidence readings before applying discontinuous
selections. Its effective premises are the prior production of
\(K\), the invertible cyclic frame, bilateral balance, the
discrete separation of memories and the archive theorem. The
register of twelve components preserves the scope of its readers;
the complete enriched history retains its other registers and
dependencies.
